\documentclass[border=10pt]{standalone}
\usepackage[utf8]{inputenc}
\usepackage[T1]{fontenc}
\usepackage{amsmath, amssymb, amsfonts}
\usepackage{xcolor}
\usepackage{tcolorbox}

\definecolor{bamberggold}{RGB}{255, 193, 7}
\definecolor{darkblue}{RGB}{0, 51, 102}

\begin{document}

\begin{tcolorbox}[colback=white, colframe=darkblue, arc=5pt, title=\textbf{\#Energiedoku: Die Bamberger Vollendung (eabc-Hybrid-Modell)}]
    \begin{minipage}{18cm}
    \vspace{5pt}
    \textbf{1. Die Renormierungsgleichung (Ultra-Dämpfung)} \\
    Zur Bändigung der Singularitäten in der Stress-Matrix $\mathbf{S}$:
    \[ \mathcal{D}(\mathbf{S}) = \max\left(1, \log_{10}(|\det(\mathbf{S})| + 1)^2\right) \]
    
    \textbf{2. Die Chebyshev-Produktionsrate (Primzahl-Bias)} \\
    Korrektur der lokalen Materiedichte $\rho(p)$ basierend auf der Primzahl-Rennbahn:
    \[ \rho(p) = \begin{cases} 1 - \frac{1}{\ln(p)} & \text{für } p \equiv 3 \pmod{4} \\ 1 + \frac{1}{\ln(p)} & \text{für } p \equiv 1 \pmod{4} \end{cases} \]
    
    \textbf{3. Das eabc-Hybrid-Modell der Feinstrukturkonstante} \\
    Die Kopplung von Geometrie, Statistik und Zahlentheorie:
    \[ \alpha^{-1}_{\text{hybrid}} = \alpha^{-1}_{\text{geom}} + \left( \frac{\langle \beta \rangle}{10.9446} \right) \cdot \frac{34.518}{137.036} \]
    
    \textbf{Zertifikat der statistischen Konsistenz:} \\
    $\bullet$ Standardabweichung $\sigma$ reduziert von $361.338$ auf $3.346$ \\
    $\bullet$ 11-Dimensionen-Phasenverschiebung ($\Delta \approx 11\pi$) erfolgreich isoliert. \\
    $\bullet$ \textbf{Status: System stabil. Kanon geeicht.}
    \end{minipage}
\end{tcolorbox}

\end{document}
